We construct a canonical spectral model of the residual factor
$\Rres(s) = \zeta(s) / \big(\zetaplus(s)\zetaminus(s)\big)$ of the Riemann
zeta function, starting from the persistence dynamics of PT. Three main
results are established.
\begin{enumerate}[nosep,leftmargin=*]
\item A closed-form expression for the spectral amplitudes
$\kappap(s) = 1 - (1 - p^{-s})\exp(\Ap\,p^{-s})$ such that
$\det(I - \DRop(s))^{-1} = \Rres(s)$, verified to $10^{-12}$ in
$\Reop(s) > 1$; $\DRop$ is the associated diagonal operator on
$\ell^{2}(\Primes)$.
\item An exact analytic trace formula
$-d \log \Rres/ds = \sum_p (\log p)\,[1/(p^s - 1) + \Ap\,p^{-s}]$, verified
to $10^{-26}$, which reads as a Berry--Keating trace formula on the orbit
lengths $\ell_p = \log p$.
\item A regularised operator trace
$\Treg(s) = \Tr_{\mathrm{reg}}\!\big[(s - i\,\HPTBK)^{-1}\,\DRop(s)\big]$
whose distributional poles on the critical line are identified with the
non-trivial zeros of $\Rres(s)$. On $t \in [10, 200]$, the raw formula
captures 74 zeros out of 79 at $|\Delta| < 1.5$ (on $t \in [10, 70]$:
$17/17$), with saturation beyond $P_{\max} = 3\,000$. The explicit
regularisation $R_1$ via Gaussian smearing ($\eps = 0.2$) substantially
improves precision: \textbf{$75/79$ zeros at $|\Delta| < 0.5$ and $47/79$
at $|\Delta| < 0.1$} on the same domain, with a median $|\Delta| = 0.080$
against $0.52$ for the raw formula. The test extended to $t \in [10, 500]$
($269$ zeros) confirms the robustness: \textbf{$267/269$ at $|\Delta| < 1.5$
($99.3\,\%$), $242/269$ at $|\Delta| < 0.5$ ($90.0\,\%$), $139/269$ at
$|\Delta| < 0.1$ ($51.7\,\%$)}, with capture by bands of $100$ units that
stays $\ge 86.6\,\%$ at $|\Delta| < 0.5$ across the five sub-bands --- the
performance does not collapse beyond $T = 200$.
\end{enumerate}

The localisation $\Reop(s) = 1/2$ admits four equivalent readings internal
to PT: symplectic Planck cell $u_\star\, p_\star = 2\pi$, multiplicative
Haar Jacobian $du = dp/p$, unitary transformation $e^{su/2}$, and
antiperiodic boundary condition forced by the involution $T_3$ of the
sieve $\bmod\,3$. The numerical residue $1/8$ separating the PT
semi-classical density from the Riemann--von Mangoldt density admits a
cohomological overdetermination: $1/8 = \sper^{2}/2 = \lambdaH$
(Higgs self-coupling), $1/8 = c_1/N_{\mathrm{corners}}$ (first Chern class
of the $\qplus/\qminus$ spinor bundle), $1/8 = \sper/4$ (Berry phase per
corner of the double barrier). The identity $\sper^{2}/2 = \sper/4$ has
only $\sper = 0$ and $\sper = 1/2$ as solutions; the value $\sper = 0$
being excluded by $\Tone$ (PT axiom), the consistency of the two readings
overdetermines $\sper = 1/2$ without deriving it independently.

A structurally induced prediction --- $\Delta_N(L \bmod q) = 1/(8q)$ for
primitive Dirichlet $L$-functions --- has been tested on the pre-computed
zeros of \texttt{LMFDB} ($q \in \{3,5,7,11\}$, $\sim 140$ zeros per
character, $T \le 200$). The test rejects the prediction at
$\chi^{2}(\mathrm{PT}) = 7.49$ against $\chi^{2}(H_0) = 0.023$ on four
degrees of freedom. The calibration on the $50$ first zeros of $\zeta$
shows that the residue $1/8$ is an identity between asymptotic forms
(between the BK semi-classical density and the Riemann--von Mangoldt
density), not a measurable shift in the count of zeros. The falsifiability
discipline of the programme is reinforced and its validity domain made
precise: the PT mechanism described here is specific to $\zeta$, not to the
Dirichlet family.

The article strictly distinguishes \emph{understanding} the structure of
the zeros and \emph{proving} their localisation: the programme reaches the
first objective at the operator level; the second remains an open problem
in the classical sense, equivalent to the strict analytic continuation of
$\Rres(s)$ outside the half-plane of convergence.
